\subsection{Ultrametric absolute values}

A mapping f of K to \(R_{+}\) is called an ultrametric absolute value if it satisfies conditions ( \(VA_{I}\) ), ( \(VA_{II}\) ) and ( \(U_{1}\) ) (which obviously implies that f is an absolute value).

PROPOSITION 3. Let \(f\) be a mapping of \(\mathbf{K}\) to \(\mathbf{R}_+\) . The following properties are equivalent:

\begin{enumerate}
\def\labelenumi{(\alph{enumi})}
\item
  \(f\) is an ultrametric absolute value.
\item
  There exists a valuation \(v\) on \(\mathbf{K}\) with values in \(\mathbf{R}\) and a real number \(a\) such that \(0 < a < 1\) and \(f = a^v\) .
\item
  \(f\) belongs to \(\mathcal{V}(\mathbf{K})\) and \(f(n,1) \leqslant 1\) for every integer \(n > 0\) .
\item
  For all \(s > 0, f^s\) is an absolute value.
\end{enumerate}

For every real number \(c\) such that \(0 < c < 1\) , the mapping \(t \mapsto c^t\) is an isomorphism of the ordered group \(\mathbf{R}\) (with the opposite ordering to the usual ordering) on the ordered group \(\mathbf{R}_+^*\) ; this shows the equivalence of (a) and (b). Clearly (a) implies (c); (c) implies (d), for we deduce from (c) that

\[
(f (n. 1)) ^ {s} \leqslant 1 \leqslant n
\]

for every integer \(n > 0\) and Proposition 2 of no. 1 shows that \(f^s\) is an absolute value. Finally (d) implies (a): for if \(f^s\) is an absolute value, it satisfies \((\mathbf{U}_2)\) and hence \(f\) satisfies \((\mathbf{U}_{2^{1 / s}})\) for all \(s > 0\) and therefore also \((\mathbf{U}_1)\) , letting \(s\) tend to \(+\infty\) .

COROLLARY. If K is a (not necessarily commutative) field of characteristic p > 0, every function on \(\mathcal{V}(\mathbf{K})\) is an ultrametric absolute value.

Every non-zero element \(z = n.1\) ( \(n\) an integer \(>0\) ) belongs to the prime subfield \(\mathbf{F}_p\) of \(K\) and hence satisfies the relation \(z^{p-1} = 1\) , which implies \(f(z) = 1\) and we may apply Proposition 3 (c).

Given a real number \(c\) such that \(0 < c < 1\) , the formulae

\[
f (x) = c ^ {v (x)}, \quad v (x) = \log_ {c} f (x)
\]

therefore establish a one-to-one correspondence between ultrametric absolute values on K and valuations on K with real values. The improper valuation corresponds to the improper absolute value (General Topology, Chapter IX, § 3, no. 2). Let \(v_{1}, v_{2}\) be two valuations on K with real values and \(f_{1}, f_{2}\) the corresponding absolute values; for \(v_{1}\) and \(v_{2}\) to be equivalent, it is necessary and sufficient that \(f_{1}\) and \(f_{2}\) be so: for to say that \(v_{1}\) and \(v_{2}\) are equivalent amounts to saying that the relations \(v_{1}(x) \geqslant 0\) and \(v_{2}(x) \geqslant 0\) are equivalent or again that the relations \(f_{1}(x) \leqslant 1\) and \(f_{2}(x) \leqslant 1\) are equivalent; it is therefore sufficient to apply Proposition 5 of General Topology, Chapter IX, §3, no. 2. Moreover (loc. cit.) for the topologies defined on K by \(f_{1}\) and \(f_{2}\) to be identical, it is necessary and sufficient that \(f_{1}\) and \(f_{2}\) be equivalent.

